\section{Related Work}
\label{sec:related}

Our work intersects two primary research areas: computational approaches to narrative quality evaluation and the emerging paradigm of LLM-based evaluation. We review representative work in each area and identify the gaps motivating our framework.

\subsection{Narrative Quality Evaluation}

The assessment of narrative quality has long challenged both literary criticism and computational approaches. Traditional natural language generation metrics including BLEU \cite{papineni2002bleu}, ROUGE \cite{lin2004rouge}, and perplexity capture surface-level textual similarity but correlate poorly with human judgments of narrative quality \cite{reiter2018structured,novikova2017we}. Embedding-based approaches like BERTScore \cite{zhang2019bertscore} improve semantic coverage through contextual representations, yet still struggle with discourse-level properties spanning extended narratives.

Recent work has developed narrative-specific evaluation frameworks addressing different quality dimensions. The HANNA benchmark \cite{chhun2022human} provides 1,056 human-annotated stories evaluated across six criteria including coherence, empathy, and engagement, with systematic comparison revealing that existing metrics capture different aspects of quality with limited agreement. NarraBench \cite{hamilton2025narrabench} presents a theory-informed taxonomy of narrative understanding tasks, surveying 78 existing benchmarks to systematize evaluation approaches. Yang and Jin \cite{yang2024makes} identify accuracy, completeness, appropriateness, insight, and consistency as the most frequently assessed quality dimensions, while noting the absence of unified frameworks capturing these dimensions simultaneously. The SCORE framework \cite{yi2025score} evaluates character consistency, emotional coherence, and plot tracking through dynamic state tracking. Complementary approaches have examined discourse coherence through entity distribution patterns \cite{barzilay2008modeling}, emotional trajectory analysis \cite{reagan2016emotional}, and fractal sentiment structures \cite{bizzoni2023fractality}, validating that computational methods can capture narratological properties at the surface and structural levels.

Despite these advances, existing frameworks share a common limitation: they assess individual quality dimensions without grounding in established literary-critical theory. Metrics for coherence, engagement, or character consistency are rarely validated against canonically recognized literature to establish that they capture properties valued in human literary criticism. As a result, evaluation risks measuring superficial textual properties while missing interpretive dimensions---cultural representation, emotional sophistication, philosophical depth---that distinguish canonical literature from competent but unremarkable prose.

\subsection{LLM-as-Judge Evaluation}

The emergence of LLMs as evaluators introduces new possibilities for scalable quality assessment. Zheng et al.~\cite{zheng2023judging} demonstrate that strong LLMs achieve over 80\% agreement with human preferences through MT-Bench and Chatbot Arena benchmarks, establishing viability for automated evaluation while identifying systematic biases including position bias, verbosity bias, and self-enhancement bias. Li et al.~\cite{li2024llms} provide a comprehensive survey noting that LLM judges correlate more highly with non-expert human judges than expert annotators, raising questions about whether such correlation reflects deep quality understanding or surface-level pattern matching. %PLACEHOLDER: additional LLM-as-judge references for multi-round or structured evaluation approaches

Most LLM-as-judge work focuses on conversational quality, instruction following, or general text generation rather than literary or narrative evaluation. The few studies addressing creative writing evaluation employ single-pass assessment without explicit theoretical grounding, leaving open the question of whether structured evaluation frameworks can improve reliability and interpretive depth. Meanwhile, the growing body of AI-generated narratives---from long-form generation systems facing systematic coherence degradation \cite{wu2025longeval} to hierarchical planning approaches \cite{gu2025rapid,wan2025cognitive}, multi-agent narrative architectures \cite{huot2024agents,ran2025bookworld}, and collaborative authoring tools \cite{yuan2022wordcraft,mirowski2023co}---produces increasingly sophisticated texts that require evaluation capable of distinguishing nuanced quality differences across multiple interpretive dimensions. Our work addresses this gap through multi-round iterative reflection, independent validation, and theory-driven prompting strategies that operationalize specific literary-critical frameworks for reliable interpretive evaluation.
